% ============================================================
%  RESUMO
%  Escreva entre os comentarios indicados.
%  Um bloco unico de texto, sem paragrafos separados.
%  Tamanho: 150 a 500 palavras (ABNT).
% ============================================================
\chapter*{RESUMO}
\addcontentsline{toc}{chapter}{RESUMO}

% --- ESCREVA SEU RESUMO AQUI ---
A geração procedural de conteúdo é uma abordagem utilizada no desenvolvimento 
de jogos digitais para produzir automaticamente elementos como cenários, mapas 
e fases. Entretanto, a utilização de diferentes estratégias de geração não 
garante que os conteúdos produzidos sejam completáveis, diversificados ou 
adequados à experiência dos jogadores. Nesse contexto, este projeto de pesquisa
tem como objetivo analisar comparativamente três estratégias de geração 
procedural de fases para jogos de plataforma 2D: geração aleatória, geração 
baseada em regras e geração por algoritmos genéticos. A pesquisa caracteriza-se 
como aplicada, exploratória, descritiva e experimental, com abordagem mista, 
envolvendo procedimentos quantitativos e qualitativos. Para sua realização, 
propõe-se o desenvolvimento de um protótipo de jogo de plataforma 2D, no qual 
as três estratégias compartilharão a mesma representação computacional das 
fases, o modelo físico e os critérios de avaliação. A completabilidade será 
investigada por meio de um procedimento automatizado baseado em busca em 
largura, enquanto a diversidade estrutural e o desempenho computacional serão 
analisados a partir de métricas quantitativas. Também está prevista uma 
avaliação exploratória com jogadores, utilizando sessões de jogo e 
questionários para investigar aspectos da experiência, como dificuldade, 
diversão, satisfação e variedade percebida. Espera-se identificar diferenças 
entre as estratégias quanto à completabilidade, à diversidade strutural, ao 
custo computacional e à experiência os participantes, contribuindo para a 
compreensão as características e limitações de diferentes bordagens de geração
procedural de fases em jogos e plataforma 2D.
% --- FIM DO RESUMO ---

\vspace{0.5cm}
\noindent\textbf{Palavras-chave:} \opalavraschave
\cleardoublepage
